\documentclass[11pt]{article}

% --- Packages (all standard; compiles with a stock pdflatex/texlive) ---------
\usepackage[utf8]{inputenc}
\usepackage[T1]{fontenc}
% Latin Modern: vector (Type 1) fonts. Without this, T1 encoding falls back
% to bitmap Type 3 EC fonts on texlive installs lacking cm-super, which
% renders fuzzy in PDF viewers.
\usepackage{lmodern}
\usepackage{microtype}
\usepackage[margin=1in]{geometry}
\usepackage{amsmath,amssymb}
\usepackage{booktabs}
\usepackage{array}
\usepackage{enumitem}
\usepackage{xcolor}
\usepackage[hidelinks]{hyperref}

\newcommand{\code}[1]{\texttt{#1}}

\title{\textbf{TraceRazor: A Heuristic Framework for Token-Efficiency\\
Auditing of LLM Agent Traces}\\[4pt]
\large A System Description and Honest Technical Report}

\author{Zulfaqar Hafez\\
\small\href{https://github.com/ZulfaqarHafez/tracerazor}{github.com/ZulfaqarHafez/tracerazor}}

\date{\today}

\begin{document}
\maketitle

\begin{abstract}
Large language model (LLM) agents routinely spend tokens on redundant
reasoning, repeated tool calls, verbose hedging, and unnecessary context
carry-over. We present \textbf{TraceRazor}, an offline, dependency-light tool
that ingests an agent execution trace and produces (i) a composite
\emph{Trace Auditing Score} (TAS) decomposed into interpretable sub-metrics
--- nine carry composite weight; five are detection-only diagnostics after
an effectiveness audit over 61 real traces demoted them ---
(ii) concrete fix suggestions with \emph{estimated} token savings,
and (iii) two optional companion pillars: consensus-based adaptive sampling
and a substitutability classifier. TraceRazor is fast (a few milliseconds for
typical traces) and runs without API keys for its core analysis.

We are deliberate about what this work does and does not establish. TAS is an
\emph{ordinal heuristic} whose composite weights are author-chosen; when we
calibrate them on real agent runs (tau-bench), no fit of the original metrics
predicts recoverable token waste (negative cross-validated $R^2$); adding
observation-accumulation features raises it into the positive range on two
independent datasets, the first signal that replicates, though the magnitude is
still modest ($\approx 0.1$). So TAS is a within-trace diagnostic, not yet a
strong predictor of cross-run savings. Reported
token ``savings'' are heuristic projections rather than measured before/after
re-runs --- and we now have direct evidence for why that distinction is
load-bearing: a \emph{measured intervention study} on a live coding agent ($24$
runs, audit$\to$apply$\to$re-run per pair, constant $12/12$ pass rate) found
that the tool's one auto-applied fix \emph{cost} a mean $5.6\%$ more tokens
(95\% bootstrap CI $[-11.4\%,\,+0.2\%]$) where the heuristic had projected a
saving. Rewriting the fix from that evidence and re-measuring yielded a
cost-neutral patch ($+0.7\%$, CI $[-8.9\%,\,+9.9\%]$). The sampling and
substitutability results remain preliminary (single-seed and synthetic,
respectively). We therefore present TraceRazor as a useful
engineering instrument and a reproducible measurement harness, not as a
validated scientific claim, and we devote a substantial Threats-to-Validity
section to its limitations. All numbers in this report are reproducible from the
public repository.
\end{abstract}

\section{Introduction}
Autonomous LLM agents are increasingly deployed in loops that interleave
reasoning steps and tool calls. Because each step consumes input and output
tokens, inefficiency compounds: an agent that re-derives the same conclusion,
retries a malformed tool call, or re-states its entire context on every turn can
cost several times more than a lean equivalent for the same outcome. Unlike raw
latency or cost dashboards, which report \emph{how much} was spent, practitioners
need to know \emph{where} tokens were wasted and \emph{what} to change.

TraceRazor targets this gap. It treats a trace as a sequence of typed steps
(reasoning, tool call, etc.) with token counts and optional tool-success flags,
and scores it along fourteen efficiency dimensions (nine composite, five
diagnostic; Section~\ref{sec:rationalisation}). The design goals are: (1)
\emph{offline by default}, the core analysis pulls in no network
dependencies; (2) \emph{interpretable}, every point of the composite maps to
a named sub-metric and a per-step annotation; and (3) \emph{actionable}, the
output includes fix patches and an ``optimal path'' diff.

\paragraph{Contributions.}
\begin{enumerate}[leftmargin=*]
  \item A decomposable token-efficiency score (TAS) over fourteen computed
        sub-metrics, of which nine carry default composite weight after an
        evidence-gated rationalisation (Section~\ref{sec:rationalisation});
        the weights are explicit and overridable (Section~\ref{sec:tas}).
  \item A heuristic fix generator and savings \emph{estimator}, plus an optimal-path
        recommendation (Section~\ref{sec:fixes}).
  \item Two optional pillars: consensus-based adaptive sampling
        (Section~\ref{sec:sampling}) and a substitutability classifier
        (Section~\ref{sec:subst}).
  \item A reproducible real-data evaluation --- 24 public function-calling and
        SWE-agent trajectories plus 13 ReAct trajectories from Hugging Face ---
        including a statistical validity audit of the fourteen sub-metrics over
        37 real traces (Section~\ref{sec:eval}), and a candid
        Threats-to-Validity analysis (Section~\ref{sec:limits}).
  \item A \emph{measured intervention study} on a live coding agent
        (Section~\ref{sec:livestudy}): the full audit$\to$apply$\to$re-run loop
        executed against real agent runs with bootstrap confidence intervals at
        constant task outcome, including a negative round that exposed a
        mispriced fix and the re-measurement after the fix was rewritten ---
        together with a transcript converter that makes any Claude~Code session
        auditable, and a token-accounting protocol that removes
        prompt-cache-warmth noise (observed $\pm22$k tokens on identical runs).
  \item Auditable runs: every report embeds a run manifest (input hash, tool
        version, backend, exact weights), a hermetic scoring mode, Ed25519
        report signing with a canonical-byte representation (any edited field
        $\to$ \textsc{tampered}), evidence bundles, and a
        \code{verify} command for third-party re-verification, plus a
        deterministic grounding-fidelity diagnostic (AGF)
        (Section~\ref{sec:metric-eval}).
\end{enumerate}

\section{Related Work}
TraceRazor draws on three lines of work. \emph{Token-budget-aware reasoning}
(e.g. TALE~\cite{tale}, SelfBudgeter~\cite{selfbudgeter}, ``Reasoning on a
Budget''~\cite{budget}) motivates the token-utilisation and value-adjusted
metrics. \emph{Reasoning-trace evaluation and step pruning}
(\cite{stepsurvey,steppruner,dualformer,efficientsurvey}) informs the
redundancy, loop, reasoning-depth, and optimal-path components. \emph{Agent
benchmarking and composite scoring} (\cite{toolchain,agenteval}) informs the
tool-sequence and aggregate-scoring design, and \emph{verbosity-bias}
studies~\cite{verbosity} inform the verbosity metrics. The adaptive-sampling
pillar is a direct application of \emph{self-consistency}~\cite{selfconsistency}:
sample multiple candidate answers and take the consensus.

A fourth line of work motivates the measured study in
Section~\ref{sec:livestudy}. \emph{Cost-aware agent evaluation}: Kapoor et
al.~\cite{agentsmatter} argue that agent benchmarks must jointly report cost
and accuracy --- optimising either alone misleads --- which is exactly the
constant-pass-rate constraint our measurement harness enforces (a token delta
on a task whose outcome flipped is not a saving). Empirical studies of agent
token spend~\cite{spendmoney,tokenomics} document both the magnitude and the
high run-to-run variance of coding-agent token consumption, which our
replicate design and bootstrap intervals are built to absorb; and
trajectory-efficiency interventions such as AgentDiet~\cite{agentdiet} and
DEPO~\cite{depo} show that trimming non-advancing content can reduce cost
without hurting task performance --- the same shape of intervention our
\code{goal\_anchor} fix attempts at the prompt level.

Relative to this literature, TraceRazor's contribution is integrative and
practical rather than a new learned method: it packages a battery of heuristics
behind a single interpretable score and an offline CLI/library. We make no claim
that the composite is a novel learned metric, nor that its weights are optimal.

\section{The TAS Framework}
\label{sec:tas}

\subsection{Trace model}
A trace is a tuple $(\textit{id}, \textit{agent}, \textit{framework},
v, S)$ where $v \in [0,1]$ is a caller-supplied \emph{task value} (task-success
quality) and $S = (s_1,\dots,s_n)$ is a step sequence. Each step $s_i$ carries a
type, free-text content, a token count $t_i$, and (for tool calls) a success
flag. Traces with $n < 5$ steps are not scored (too little signal). Textual
similarity $\mathrm{sim}(\cdot,\cdot)\in[0,1]$ between step contents is provided
by an injected backend; the default is a bag-of-words cosine, with an optional
embedding backend.

\subsection{Sub-metrics}
\label{sec:metrics}
Each sub-metric maps a trace to $[0,1]$ (or a percentage), higher being more
efficient. Table~\ref{tab:metrics} lists all fourteen with their default
composite weight; five carry weight zero by default --- they are computed,
reported and fix-driving, but a real-trace effectiveness audit
(Section~\ref{sec:rationalisation}) found them unable to earn a share of the
composite. Representative definitions:

\begin{itemize}[leftmargin=*]
  \item \textbf{SRR (Step Redundancy Rate).} Fraction of steps that are
        near-duplicates of a recent prior step
        ($\mathrm{sim} \ge 0.65$ for the bag-of-words backend), scanned over a
        bounded look-back window of the last $256$ steps to keep the cost
        $O(n\cdot 256)$ rather than $O(n^2)$.
  \item \textbf{LDI (Loop Detection Index).} Degree to which the
        reasoning/tool-call sequence revisits previously seen cycles. Loops are
        detected both by exact state hashing (tool name $+$ serialised
        parameters) and, for command-style tools, by a \emph{parametric}
        detector that abstracts argument literals into a command skeleton so
        that the same command template repeated for different arguments (a loop
        shape common in tool-using agents) is caught (Section~\ref{sec:hf-agentinstruct}).
  \item \textbf{TCA (Tool Call Accuracy).} Fraction of tool calls that
        succeeded on first attempt.
  \item \textbf{ISR (Information Sufficiency Rate).} Fraction of steps that add
        novelty, where a step's information gain is
        $1 - \max_{j\in W(i)} \mathrm{sim}(s_i, s_j)$ over the bounded prior
        window $W(i)$; steps below a novelty threshold are flagged.
  \item \textbf{TUR, CCE.} Token-utilisation ratio and context carry-over
        efficiency, penalising re-stated context.
  \item \textbf{RDA (Reasoning Depth Appropriateness).} Closeness of the actual
        step count to an expected count inferred from task complexity (tool
        surface area, query patterns, historical median).
  \item \textbf{DBO (Decision Branch Optimality).} Tool-sequence efficiency
        against accumulated historical sequences; with no history it falls back
        to a single-trace proxy whose retry and thrash signals key on the
        \emph{invocation} (tool $+$ parameters), not the bare tool name, so
        single-tool command agents are not structurally penalised
        (Section~\ref{sec:hf-agentinstruct}).
  \item \textbf{Verbosity: VDI, SHL, CCR.} Verbosity-density, sycophancy/hedging
        level, and a compression-ratio proxy.
  \item \textbf{GAR, CSD.} Goal-advancement ratio and cross-step semantic drift.
        Both score reasoning-bearing steps against the goal (GAR) or each other
        (CSD); the scored set includes ReAct tool-call turns whose content embeds
        a substantive thought, not just steps typed as reasoning, so the metrics
        do not collapse on tool-using agents whose reasoning is fused with action.
        Fenced code inside a scored turn is reduced to its argument literals
        (paths, quoted strings, numbers) before similarity: the literals anchor
        the step to the task while command syntax no longer dilutes the overlap
        (Section~\ref{sec:hf-agentinstruct}).
\end{itemize}

\begin{table}[t]
\centering
\caption{The TAS sub-metrics and their default weights. Non-zero weights sum
to $0.92$ and are normalised by their sum at scoring time, so the rightmost
column (\,weight\,$/0.92$) is the actual share of the composite. The five
zero-weight metrics are detection-only diagnostics, demoted by the
effectiveness audit of Section~\ref{sec:rationalisation}; any of them can be
re-enabled via a weights file. All weights are author-chosen heuristics
except OBS, which was added after it predicted real recoverable waste across
two datasets (Section~\ref{sec:betterfeatures}); the rest remain uncalibrated
(Section~\ref{sec:limits}).}
\label{tab:metrics}
\begin{tabular}{llrr}
\toprule
Code & Metric & Weight & Share \\
\midrule
SRR & Step Redundancy Rate            & 0.17 & 18.5\% \\
LDI & Loop Detection Index            & 0.13 & 14.1\% \\
TCA & Tool Call Accuracy              & 0.13 & 14.1\% \\
RDA & Reasoning Depth Appropriateness & 0.10 & 10.9\% \\
ISR & Information Sufficiency Rate     & 0.10 & 10.9\% \\
TUR & Token Utilisation Ratio         & 0.10 & 10.9\% \\
CCE & Context Carry-over Efficiency   & 0.10 & 10.9\% \\
OBS & Observation Token Share         & 0.06 & \phantom{0}6.5\% \\
CCR & Caveman Compression Ratio       & 0.03 & \phantom{0}3.3\% \\
\midrule
DBO & Decision Branch Optimality      & 0 & diagnostic \\
VDI & Verbosity Density Index         & 0 & diagnostic \\
SHL & Sycophancy/Hedging Level        & 0 & diagnostic \\
GAR & Goal Advancement Ratio          & 0 & diagnostic \\
CSD & Cross-Step Semantic Drift       & 0 & diagnostic \\
\midrule
    & \textbf{Total (composite)}      & \textbf{0.92} & 100\% \\
\bottomrule
\end{tabular}
\end{table}

\subsection{Composite score}
Let $m_k\in[0,100]$ be the $k$-th sub-metric and $w_k$ its weight. The raw
structural score is the normalised weighted mean
\begin{equation}
\textit{raw} \;=\; \frac{\sum_k w_k\, m_k}{\sum_k w_k},
\qquad \sum_k w_k = 0.92
\;\text{(non-zero default weights; Table~\ref{tab:metrics})} .
\end{equation}
The reported TAS applies a \emph{Task Value Integration} (TVI) multiplier so that
a structurally clean trace that nonetheless \emph{failed} its task is penalised:
\begin{equation}
\textit{TAS} \;=\; \textit{raw}\times\bigl(0.7 + 0.3\,v\bigr).
\label{eq:tvi}
\end{equation}
At $v=1$ the multiplier is $1$ (no change); at $v=0$ it caps the score at $0.7$
of raw. The constants $0.7/0.3$ are a design choice, not a fitted parameter.

\subsection{Verbosity alert}
An aggregate verbosity score
$\textit{AVS} = 0.45(1-\textit{VDI}) + 0.30\,\textit{SHL} + 0.25\,\textit{CCR}$
is reported alongside TAS; $\textit{AVS} > 0.40$ raises a verbosity alert. As
with the composite weights, these coefficients and threshold are heuristic.

\subsection{Calibrating the weights}
\label{sec:calib}
Because the default weights are heuristic, we provide a calibration procedure
that fits them to data for a stated objective (here, \emph{recoverable token
waste}). Given a dataset of traces each labelled with a measured recoverable
fraction $r\in[0,1]$ (from before/after re-runs at constant task quality), we set
the target efficiency $y = 1-r$ and the feature vector $x$ to the
normalised sub-metrics, then solve the simplex-constrained least squares
\begin{equation}
\min_{\mathbf{w}}\; \lVert X\mathbf{w} - \mathbf{y}\rVert_2^2 + \lambda\lVert \mathbf{w}-\tfrac{1}{13}\mathbf{1}\rVert_2^2
\quad\text{s.t.}\quad \mathbf{w}\ge 0,\ \textstyle\sum_k w_k = 1 ,
\end{equation}
where the constraint $\sum_k w_k = 1$ makes the engine's normalised composite
equal $X\mathbf{w}$ exactly, and $\lambda$ optionally regularises toward uniform
weights. The tool reports in-sample $R^2$, $k$-fold cross-validated $R^2$, and
the default-weights baseline, and emits a weights file the engine loads at
runtime. Measured before/after pairs may be supplied as raw, LangSmith, or OTEL
traces through an import adapter; token totals are read from the auditor so the
two sides are comparable across formats.

We validate this procedure directly on real agent runs rather than on synthetic
data; the next two subsections report that real-data calibration and the feature
work it motivated. The engine default weights are left heuristic (except OBS,
which earns its place on the real-data evidence below) until a user calibrates on
their own measured runs.

\subsection{Real-data calibration on \texorpdfstring{$\tau$}{tau}-bench (a negative result)}
\label{sec:realcalib}
We then calibrated on real agent runs. tau-bench~\cite{taubench} publishes
trajectories for GPT-4o and Claude~3.5 Sonnet on the airline and retail domains;
of $1{,}980$ episodes, $1{,}183$ succeed (reward $=1$). Because the same task is
solved by multiple models and repeated trials, successful runs of one task at
differing token cost form before/after pairs \emph{at equal quality} (both
completed the task). Restricting to same-model, different-trial pairs (the purest
within-agent signal) yields $233$ pairs; allowing cross-model pairs yields $147$.

The result is negative and we report it as such. Table~\ref{tab:realcalib} shows
that on real data the calibrated weights reach a \emph{negative} cross-validated
$R^2$ ($-0.11$), i.e.\ no convex combination of the thirteen sub-metrics predicts
real recoverable token waste better than predicting the mean. Every individual
sub-metric correlates weakly with the target ($|r|<0.3$; the strongest is CCE at
$-0.29$). In other words, the original metrics genuinely do not track what drives
recoverable token differences on real runs.

\begin{table}[t]
\centering
\caption{Calibration on real tau-bench before/after pairs (5-fold CV). Negative
$R^2$ means the fit does not beat predicting the mean; TAS's metrics do not
predict real recoverable waste here.}
\label{tab:realcalib}
\begin{tabular}{lrr}
\toprule
Weights & Train $R^2$ & CV $R^2$ \\
\midrule
Default (heuristic)    & --     & $-0.37$ \\
Calibrated (best fit)  & $-0.09$ & $-0.11$ \\
\bottomrule
\end{tabular}
\end{table}

\paragraph{Interpretation.} The likely cause is a target/score mismatch: the
cross-run token delta between a verbose and a lean successful run is driven
mostly by how many steps the agent took and which path it chose, whereas TAS
measures \emph{within-trace} structural redundancy. A run can be internally
non-redundant (good TAS) yet still take a longer valid path than a leaner run
(high recoverable fraction), which breaks the correlation. The honest conclusion
is that TAS is a within-trace diagnostic, useful for flagging redundancy inside a
single run, but it is \emph{not}, as currently defined, a validated predictor of
how many tokens are recoverable relative to a leaner run. Closing that gap needs
richer or non-linear features (e.g.\ path-level and step-count modelling) and a
target measured on within-agent optimisation rather than cross-run comparison.

\subsection{Better features: observation-accumulation signals}
\label{sec:betterfeatures}
Acting on that interpretation, we added candidate features that measure
\emph{context accumulation} rather than reasoning redundancy: the tool-observation
token share, output compressibility, repeated observations, and redundant tool
calls (all computed from fields the parsers already capture, emitted as
diagnostics and not in the composite). We then tested whether they raise
predictiveness on real data, across two independent datasets:
tau-bench and tau2-bench (the latter adds the telecom domain and the
Claude~3.7~Sonnet / GPT-4.1 models). Table~\ref{tab:features} shows the features
raise the cross-validated $R^2$ over metrics-alone on both datasets, and the
combined model is positive on both ($+0.08$ and $+0.12$) where metrics alone were
negative or near-zero. The one consistent driver is the observation token share
(correlation $+0.28$ and $+0.33$ with efficiency). This is the first signal here
that predicts real recoverable waste and replicates across datasets.

\begin{table}[t]
\centering
\caption{Adding observation-accumulation features raises cross-validated $R^2$
against real recoverable waste on two independent datasets (5-fold CV).}
\label{tab:features}
\begin{tabular}{lrrr}
\toprule
Dataset (pairs) & default & metrics only & metrics + features \\
\midrule
tau-bench (233)    & $-0.37$ & $-0.11$ & $\mathbf{+0.08}$ \\
tau2-bench (1055)  & $-0.15$ & $+0.01$ & $\mathbf{+0.12}$ \\
\bottomrule
\end{tabular}
\end{table}

\paragraph{Promotion and honest magnitude.} Having cleared the evidence-gated bar
on both datasets, the consistent driver, observation token share, is promoted
into the composite as the OBS metric (Table~\ref{tab:metrics}, weight $0.06$,
$4.8\%$ share); the remaining candidate features (stale-observation retention,
context growth) stay diagnostic in the report. The absolute $R^2$ is still modest
(about $0.1$): OBS improves the score's grounding and crosses the ``better than
the mean'' bar consistently, but it does not yet make TAS a strong predictor of
recoverable waste. We keep OBS at a deliberately modest weight, and note its
positive direction was validated on conversational tool-agent datasets and should
be re-checked on other domains (e.g.\ coding agents) via the calibration tool.

\paragraph{Where the ceiling is.} We pushed further: we added path/length
features (step count, mean tokens per step, longest reasoning run, state-revisit
rate, tool diversity) and tested non-linear models. None raised cross-validated
$R^2$ above the current level on either dataset --- the extra features are largely
collinear with the existing metrics plus OBS, and on tau2-bench ($n{=}1055$) the
regularised convex fit ($R^2 = 0.12$) beat both ridge regression ($0.09$) and
gradient boosting ($0.04$), so there is no hidden non-linear signal these
features are missing. The structural-feature ceiling for predicting cross-run
recoverable waste is therefore about $0.1$. Materially exceeding it would need
fundamentally richer signal (semantic/path-embedding features, per-token
attribution), which is a research effort rather than a feature add, and remains
future work.

\section{Fix Generation and Savings Estimation}
\label{sec:fixes}
For each detected inefficiency, TraceRazor emits a typed fix patch
(e.g. \emph{context compression}, \emph{termination guard},
\emph{tool-schema tightening}) with an \emph{estimated} per-run token saving
computed from per-fix heuristics (for example, projecting that a fraction of a
flagged step's tokens is recoverable). It also produces an \emph{optimal path}:
a keep/trim/delete annotation over the steps with a projected token total.

\paragraph{Estimates are not measurements.} The headline ``tokens saved'' and
``\$/month'' figures are the \emph{sum of these per-fix projections}. They are
not the result of applying the fixes and re-running the agent. TraceRazor
provides a \code{bench} command precisely so that a user can supply a real
``after'' trace and obtain a \emph{measured} before/after delta; only that path
yields a measured saving. This distinction is load-bearing and is restated in
Section~\ref{sec:limits}; Section~\ref{sec:livestudy} quantifies it directly:
on a live coding-agent workload the projection for the one auto-applied fix had
the wrong \emph{sign} (estimate accuracy $-102\%$), and the fix's patch text
was rewritten from that measurement. The current \code{goal\_anchor} wording
instructs the agent to keep the objective as working context and skip
non-advancing actions \emph{without} restating the objective --- the earlier
wording's per-step restating ritual was the measured cost driver.

\section{Pillar 2: Adaptive Sampling}
\label{sec:sampling}
The sampling pillar applies self-consistency~\cite{selfconsistency} to agent
final answers: sample $K$ candidates, canonicalise them (numeric/whitespace/
key-order normalisation), and select by exact-match plurality vote. An
\emph{adaptive-$K$} variant reduces $K$ when recent steps reach consensus quickly
and restores it otherwise, trading a small accuracy cost for fewer tokens. The
in-package consensus selector uses \textbf{honest majority vote with a
deterministic tiebreak and no access to ground-truth answers}.

We emphasise two caveats up front (expanded in Section~\ref{sec:limits}): the
benchmark below is \textbf{single-seed over 50 tasks with no confidence
intervals}, and the adaptive rule conflates consensus with correctness. The
``Pareto'' framing is therefore \emph{directional}, not a significance claim.

\section{Pillar 3: Substitutability Classifier}
\label{sec:subst}
This pillar predicts whether a cached response can substitute for a fresh LLM
call, using embedding-similarity features and an sklearn classifier. We report it
for completeness but stress that the shipped evaluation is a \textbf{pipeline
smoke test on synthetic templated pairs}, where the binary label is essentially
``same intent.'' Under this construction the task collapses to topic similarity,
which the embedding feature trivially recovers; the resulting near-perfect AUC
reflects \emph{label leakage by design} and must not be read as a generalisation
estimate. We include it to document the pipeline, not to make a performance
claim.

\section{Evaluation}
\label{sec:eval}
All results in this section are reproducible from the repository
(Section~\ref{sec:repro}).

\subsection{Audits on real public agent trajectories}
\label{sec:public-audits}
We converted 24 real public agent runs (tau-bench airline/retail with GPT-4o and
Claude~3.5 Sonnet; SWE-agent edit-format variants) into the trace schema and
audited them hermetically (fresh state per audit, so numbers are independent of
audit order). Table~\ref{tab:taubench} reports mean TAS and mean step-redundancy
per model/domain under the current scorer (the nine-signal composite of
Section~\ref{sec:rationalisation}); the corpus mean TAS is $78.4$.
Redundancy here reflects the \emph{responsiveness} rules adjudicated in on real
traces (Section~\ref{sec:metric-eval}): a step answering a new user turn, a
successful retry of a failed call, and a verification re-run after a state
change are never counted as redundant.
Because token counts for some sources are approximated from text length,
absolute token figures should be read as estimates; the \emph{relative} ordering
is the informative part. The ordering is stable and interpretable: retail runs
score highest ($88$--$93$), GPT-4o airline runs lowest ($59$), tracking the
step-redundancy column. For SWE-agent, the compact \code{xml} edit format used
about \textbf{52\% fewer tokens} than the verbose \code{cursors} format
($3{,}636$ vs $7{,}553$ tokens) at comparable TAS ($\approx 72$--$76$). The
audits are regenerated by a script (\code{benchmark/run\_benchmarks.py}) so
repository numbers track the current scorer; we do not ship a synthetic
benchmark.

\begin{table}[t]
\centering
\caption{Audit summary over real public trajectories (means, hermetic audits,
current scorer). TAS is the ordinal composite; SRR is step-redundancy (lower
is better).}
\label{tab:taubench}
\begin{tabular}{llrr}
\toprule
Model & Domain & Mean TAS & Mean SRR \\
\midrule
GPT-4o          & airline & 59.3 & 15.5\% \\
GPT-4o          & retail  & 93.1 & 6.1\% \\
Claude 3.5 Sonnet & airline & 74.2 & 24.0\% \\
Claude 3.5 Sonnet & retail  & 87.7 & 3.7\% \\
\bottomrule
\end{tabular}
\end{table}

\subsection{Real ReAct trajectories from Hugging Face (AgentInstruct)}
\label{sec:hf-agentinstruct}
The $\tau$-bench traces above are function-calling assistants. To test the
auditor against a different real agent style --- tool-using \emph{ReAct} agents
that interleave a natural-language thought with a shell or SQL action --- we
sourced trajectories from the Hugging Face dataset \code{zai-org/AgentInstruct}
(formerly \code{THUDM/AgentInstruct})~\cite{agentinstruct}, a corpus of
$1{,}866$ ReAct interactions across six task domains. We converted a sample
spanning the operating-system (bash) and database (SQL) splits into the trace
schema (\code{tools/convert\_agentinstruct.py}); the resulting corpus of
thirteen traces is audited end-to-end by a \code{cargo test} statistics gate
(\code{crates/tracerazor-cli/tests/huggingface\_real\_data.rs}) and summarised by
\code{benchmark/hf\_audit\_stats.py}, with every audit run in a fresh state
directory so measurements are independent of audit order. Provenance and a live
dataset-viewer fetch path are documented in the repository.

\paragraph{Few-shot scaffolding is not agent behaviour.} Each upstream row
embeds the dataset's fixed one-shot demonstration (the same ``count the files
in \code{/etc}'' exchange in every \code{os} row) before the real task, marked
by a \code{loss{=}false} flag on its assistant turns. Our first conversion
audited those turns as agent steps, which pseudo-replicated one canned
trajectory into every trace and mis-anchored the goal metrics (every trace
``began'' with a step about \code{/etc} regardless of its task). Excluding
scaffolding by the loss flag (with a textual-boundary fallback) changed the
statistics materially: mean TAS fell from $82.8$ to $78.0$ (under the
then-current fourteen-metric composite) --- the demo steps
had been \emph{padding} every score --- while goal alignment rose
($\mathrm{GAR_{norm}}$ $0.26\to0.35$) once the metrics no longer fought the
demo/task mismatch. We flag this as a general hazard when auditing
conversation-formatted agent corpora.

Auditing the corpus drove four data-driven fixes (composite weights unchanged):

\paragraph{Structural blind spots: parametric loops and single-tool DBO.} LDI
keyed loops on an exact state hash and never fired ($\mathrm{LDI_{norm}}=1.0$
everywhere) even on \code{os\_6}, which runs the same template once per file
(\code{grep -o "Linux" <FILE> | wc -l}, four times); the parametric detector
(Section~\ref{sec:metrics}) abstracts argument literals into a command skeleton
and flags $\ge 3$ repeats, scoped to command-style tools. On the
de-contaminated corpus \code{os\_6} scores $\mathrm{LDI_{norm}}=0.333$ with a
loop-termination fix, while progressive pipelines (\code{os\_0}) are correctly
not flagged. Similarly, DBO's cold-start fallback keyed retries on the bare
tool name, capping a bash operator --- one tool by construction --- near the
$0.5$ floor regardless of conduct; keying on the \emph{invocation} (tool $+$
parameters) moved corpus mean $\mathrm{DBO_{norm}}$ $0.59\to0.88$, with the one
genuine-failure trace (\code{os\_5}) now the only one below the ceiling.

\paragraph{Semantic blind spot: reasoning fused into tool calls.} GAR and CSD
scored only reasoning-typed steps, but a ReAct agent's goal-relevant reasoning
lives inside the tool-call turn (``Think: \dots\ Act: bash \dots''), so both
collapsed on tool agents ($\mathrm{GAR_{norm}}$ as low as $0.0$). Two changes
fixed this: tool-call steps carrying substantive reasoning prose ($\ge 12$
words) are now scored, and fenced code inside a scored turn is reduced to its
argument literals. The second is data-driven in a specific way: an ablation
that stripped fences wholesale made scores \emph{worse} (corpus
$\mathrm{CSD_{norm}}$ $0.415\to0.353$) because the code's paths and quoted
strings are exactly the tokens shared with the natural-language goal --- so the
shipped reduction keeps the literals and drops only command syntax, the inverse
of the LDI skeleton. $\mathrm{GAR_{norm}}$ averages $0.348$ on the
de-contaminated corpus, up from $0.202$ at the start of the exercise.

\paragraph{Most real ReAct task runs are shorter than the analysis floor.}
With scaffolding excluded, $9$ of $13$ traces ($\approx69\%$) fall below the
default five-step analysis floor --- real task runs routinely finish in three
or four steps. Rather than silently excluding that majority, the CLI gained a
\code{-{}-min-steps} opt-in (default unchanged; clamped to $\ge 2$, where
pair-based metrics carry less evidence): with \code{-{}-min-steps 2} the gate
audits $13/13$ traces (mean TAS $86.1$, $20$ fix patches), so the floor's
coverage cost on real data is now measured and recoverable instead of assumed.

Table~\ref{tab:agentinstruct} reports the corpus statistics after the fixes.
The value here is not the headline score but the corrected blind spots and the
scaffolding hazard, which neither a synthetic benchmark nor a
function-calling-only evaluation would have surfaced --- concrete evidence for
testing heuristic auditors against \emph{diverse} real agent styles.

\begin{table}[t]
\centering
\caption{TraceRazor on the real Hugging Face AgentInstruct corpus after
few-shot-scaffolding exclusion (13 traces; 4 analysable at the default 5-step
floor). Normalised metric means ($1.0=$ no waste detected).}
\label{tab:agentinstruct}
\begin{tabular}{lrl}
\toprule
Quantity & Value & Note \\
\midrule
Mean TAS (default floor) & 79.8  & all ``Good''; median 80.3 \\
Mean SRR$_\text{norm}$ & 0.650 & step redundancy detected \\
Mean LDI$_\text{norm}$ & 0.833 & was $1.000$; \code{os\_6} parametric loop at $0.333$ \\
Mean GAR$_\text{norm}$ & 0.348 & was $0.202$; reasoning-aware $+$ literal-anchored \\
Mean CSD$_\text{norm}$ & 0.487 & full reasoning flow (was a degenerate pair) \\
Mean DBO$_\text{norm}$ & 0.875 & was $\approx0.57$; invocation-keyed cold start \\
Mean OBS$_\text{norm}$ & 0.930 & high tool-I/O share \\
Fix patches (sum)     & 11    & across the four analysable traces \\
Full corpus (\code{-{}-min-steps 2}) & 13/13 & mean TAS $86.1$; $20$ fix patches \\
\bottomrule
\end{tabular}
\end{table}

\subsection{A measured intervention study on a live coding agent}
\label{sec:livestudy}
Every evaluation above audits traces; none measures whether \emph{acting on an
audit} actually saves tokens. This subsection reports the first closed-loop
measurement: the full audit$\to$apply$\to$re-run protocol executed against a
live agent, with the product's own emitted fix as the only manipulated
variable. Following~\cite{agentsmatter}, cost is never reported without task
outcome: the harness refuses to call a token delta a ``saving'' on any task
whose pass flag flips between conditions, and exits non-zero.

\paragraph{Design.} The subject agent is Claude~Code (v2.1.170, headless,
model Haiku~4.5) --- a third real agent style after the function-calling and
ReAct corpora: an interactive coding agent. The workload is six small Python
engineering tasks (fix an off-by-one, implement to a failing spec, rename an
API across files, deduplicate drifted helpers, implement a CSV filter, repair
package imports), each with an objective binary outcome: the task's pytest
suite, failing at start, must pass when the agent finishes. Each task runs
twice per condition (12 pairs), because run-to-run variance in coding-agent
token spend is large~\cite{spendmoney}. The tool envelope (pytest-only shell,
file read/edit, 25-turn cap) is identical across conditions. Per pair:
(1)~run the stock agent, convert the session transcript to a trace, record the
pytest outcome as task value; (2)~\code{audit -{}-hermetic}, then \code{apply}
the safe-fix subset to a prompt file (the product's default behaviour);
(3)~re-run the same task with that file appended to the system prompt --- the
only delta; (4)~measure with \code{bench} per pair and aggregate with a
seeded $10{,}000$-resample percentile bootstrap. Total API cost for the study
was $\approx\$1.30$; all traces, fixes, prompts and run logs are committed.

\paragraph{Token accounting.} Converting interactive-agent transcripts
surfaced a measurement hazard worth reporting on its own: the first turn's
prompt-prefix encoding ($\approx20$k tokens of harness system prompt) is
billed as a cache \emph{write} or a cache \emph{read} depending on cache
warmth at launch --- we observed $\pm22$k-token swings between otherwise
identical runs, several times the size of the behavioural signal. The
converter therefore uses \emph{marginal} accounting: per turn,
$\mathrm{input} + \mathrm{cache\ creation} + \mathrm{output}$, excluding cache
reads and the first turn's prefix encoding (constant across compared
conditions by construction); the convention is stamped into every trace's
metadata. Input-side accumulation dominating agent cost is consistent
with~\cite{tokenomics}.

\paragraph{Round 1: the shipped fix, measured (a negative result).} On all 12
stock runs the audits emitted exactly one auto-applicable fix,
\code{goal\_anchor}, with projected savings of $0.5$--$2.4$k tokens/run.
Measured (Table~\ref{tab:livestudy}): mean token reduction
$\mathbf{-5.6\%}$ (95\% CI $[-11.4\%,\,+0.2\%]$) at constant $12/12$ pass rate
--- the patch \emph{added} cost. Estimate accuracy
(measured\,/\,projected savings) was $-102\%$: the projection had the wrong
sign. The mechanism is visible in the patch text: the then-current wording
instructed the agent to ``restate the stated task objective in one sentence
before each reasoning step,'' a per-turn standing cost; and because goal
advancement (GAR) passed on every before-trace, there was no off-path
exploration for the anchor to recover. The remediation was triggered by the
trajectory-path-entropy detector, whose conversational prior misreads
read$\to$edit$\to$test cycles as drift (Section~\ref{sec:limits}).

\paragraph{The intervention, revised from the measurement.} The patch wording
was rewritten --- anchor silently, skip non-advancing actions, ``do not
restate the objective or summarise progress unless explicitly asked'' ---
leaving detection and the conservative estimate unchanged. This is the same
intervention shape that trajectory-reduction methods validate at the
context level~\cite{agentdiet,depo}: remove non-advancing content, never add
ritual.

\paragraph{Round 2: the revised fix, re-measured.} Fresh audits and 12 fresh
after-runs: mean token reduction $\mathbf{+0.7\%}$ (95\% CI
$[-8.9\%,\,+9.9\%]$), mean TAS delta $-0.0$, pass rate $12/12 \to 12/12$. The
standing cost is gone, and we are explicit that \emph{zero is the correct
value for this workload}: a drift remediation applied to non-drifting runs
should do nothing. Whether the anchor produces positive savings on genuinely
drifting agents (GAR failing) is the designed next experiment. Per-pair deltas
span $-30\%$ to $+24\%$ between identical configurations, so no single pair is
evidence of anything; the aggregate interval is the unit of claim.

\begin{table}[t]
\centering
\caption{The measured intervention study: 6 tasks $\times$ 2 replicates per
round, Claude~Code (Haiku~4.5), seeded 10{,}000-resample bootstrap 95\% CIs,
pass rate constant at $12/12 \to 12/12$ in both rounds. Negative ``reduction''
is an added cost. Round~1 measured the shipped \code{goal\_anchor} patch;
round~2 the rewrite derived from round~1's diagnosis.}
\label{tab:livestudy}
\begin{tabular}{lrrr}
\toprule
Round & Mean token reduction (95\% CI) & Mean TAS $\Delta$ & Estimate accuracy \\
\midrule
1 (restate-each-step)     & $-5.6\%\ [-11.4,\,+0.2]$ & $-1.5$ & $-102\%$ \\
2 (anchor-without-restate) & $+0.7\%\ [-8.9,\,+9.9]$  & $-0.0$ & $\approx 0\%$ (wide CI) \\
\bottomrule
\end{tabular}
\end{table}

\paragraph{What this establishes --- and what it does not.} It establishes
that the measurement harness works as designed: it caught a mispriced fix
(ours), quantified the gap between projection and outcome, and verified the
repair, all at constant task success. It does \emph{not} establish that
TraceRazor saves tokens on coding agents --- on this on-track workload there
was nothing to save, and only one fix type fired. The study's value is
methodological: the loop from heuristic projection to measured outcome is now
closed, cheap ($\approx\$1$ per round), and committed as a reproducible kit
(\code{benchmark/live/}).

\subsection{Metric validity, grounding fidelity, and verifiable runs}
\label{sec:metric-eval}
A statistical audit of the fourteen sub-metrics over $37$ real traces
($20$ $\tau$-bench, $13$ AgentInstruct, $4$ SWE-agent; independent state per
audit) quantified how the composite behaves in practice. Three findings stand
out. First, the task-value multiplier --- not the structural metrics --- drives
the final score: $r(\mathrm{TAS},\allowbreak \mathrm{task\_value\_score})=0.89$,
while the structural metrics separate failed from passing $\tau$-bench runs by
only $4.9$ raw points. Second, realised influence diverges from the documented
weights because influence scales with each metric's empirical variance:
TUR (nominal $7.9\%$) carries $28\%$ of raw-TAS variance ($r=0.89$), whereas
the verbosity trio VDI/SHL/CCR ($12.7\%$ nominal) contributes $\approx 0$, with
VDI mildly \emph{anti}-correlated. Third, normalisation ranges are
miscalibrated at both ends: GAR never exceeds $0.62$ and CSD never exceeds
$0.68$ on any real trace (a near-constant drag rather than a discriminator),
while LDI sits at its $1.0$ ceiling on $78\%$ of traces; DBO and TCA are
strongly collinear ($r=0.81$). These findings are no longer just noted: they
drove the metric rationalisation of Section~\ref{sec:rationalisation}, which
re-derived them on a larger corpus and removed the offending metrics from the
default composite. A subsequent precision pass (driven by
step-level reviewer adjudication of two traces) added \emph{responsiveness}
exemptions to SRR and LDI --- new-input, fail$\to$retry, and
verification-after-mutation --- and a guard that a successful state-changing
call is never a delete candidate; on the adjudicated airline trace this took
the delete-recommendation precision from $1/6$ correct to $6/6$, and shifted
corpus SRR distributions accordingly (airline mean $35.9\%\to15.5\%$).

\paragraph{AGF: a grounding-fidelity diagnostic for auditable runs.} Guided by
the trajectory-hallucination literature~\cite{trajhall}, the report now carries
\textbf{Action/Claim Grounding Fidelity}: the share of tool-call argument
literals (paths, quoted strings, numbers) that appear in prior context, and of
final-answer literals that appear in \emph{environment-provided} text (task,
observations). Matching is deterministic literal containment --- no model, no
similarity backend --- so the diagnostic itself is reproducible byte-for-byte,
and every failure is itemised (\code{ungrounded[]}) with the step it occurred
in. Derived values (arithmetic the agent performed) are deliberately counted as
ungrounded: they are precisely the claims an auditor must verify by hand.
Reviewer adjudication exposed the first tokenizer's failure mode --- markdown
emphasis, shell variables, regex/awk/glob patterns, and code the agent itself
writes were flagged as ``ungrounded claims'' --- so the extractor now strips
syntax classes, treats content-creation params (edits) as artifacts rather
than assertions, and counts a step's own \code{input\_context} as evidence.
A corpus-wide acceptance test asserts zero syntax artifacts in
\code{ungrounded[]}. After the rewrite AGF averages $0.951$ on the
AgentInstruct corpus and $0.966$ on the function-calling/SWE corpus, with the
failure/retry trace lowest ($0.875$).
AGF is reported as a diagnostic, not folded into TAS: the variance analysis
above shows what uncalibrated weights do to realised influence.

\paragraph{Run manifests and third-party verification.} A score is only
evidence if a third party can check what was scored, by what, and under which
configuration~\cite{provagent}. Every audit now embeds a \emph{run manifest}:
SHA-256 of the raw input trace bytes, tool version, RFC-3339 timestamp, the
similarity backend that \emph{actually} ran (silent embedding$\to$BoW fallbacks
become recorded fact), the exact composite weights (inline, plus their
SHA-256), the step floor, and the store-derived baselines that influenced
scoring, if any. A new \code{-{}-hermetic} mode reads nothing from and writes
nothing to local state, making the score a pure function of (trace, config,
version) --- and \code{tracerazor verify <report> <trace>} re-checks the hash,
then re-scores hermetic bag-of-words runs from the manifest's configuration
and compares every normalised metric exactly, exiting non-zero on a single
flipped byte or score divergence. Non-reproducible conditions (embedding
backends, store-influenced baselines) are detected and reported as
hash-and-integrity-only verification rather than silently glossed over.

Reports can additionally be \emph{signed}: an Ed25519 signature over a
canonical byte representation of the report (timing fields zeroed, signature
fields nulled, floating-point values normalised through a JSON round-trip so
sign-time and verify-time bytes agree exactly). \code{verify} then proves both
provenance and integrity --- editing any field of a signed report, from the
headline score to a single per-step annotation, yields a \textsc{tampered}
verdict and a non-zero exit; unsigned reports verify as
``unsigned'' rather than being conflated with verified ones. An evidence
bundle (\code{export -{}-bundle}) packages trace, report, weights and
checksums into a single verifiable zip for handoff to a third party. Getting
canonical bytes stable required fixing two real nondeterminism sources ---
\code{HashMap}-ordered loop groups in LDI and last-digit drift in f64 JSON
round-trips --- both of which would otherwise have made honest reports fail
verification intermittently.

\subsection{Metric rationalisation: an evidence-gated composite}
\label{sec:rationalisation}
Acting on the validity audit above, we re-derived its findings on a larger
corpus and pruned the composite accordingly. The evaluation
(\code{benchmark/metric\_effectiveness.py}, output committed as
\code{docs/metric\_effectiveness.md}) audits \textbf{61 real traces} --- the
24 public function-calling/SWE trajectories, the 13 AgentInstruct ReAct
trajectories, and the 24 live coding-agent traces from
Section~\ref{sec:livestudy} --- and applies three criteria stated
\emph{before} the data is read: a metric keeps composite weight only if it is
(C1) \emph{discriminative} (corpus sd $\ge 0.05$), (C2) \emph{range-sane}
(corpus max $\ge 0.80$; a normalised metric that never approaches its top on
real data is a constant drag, not a discriminator), and (C3)
\emph{non-redundant} ($|r| < 0.85$ against every kept metric).

Five metrics failed: \textbf{GAR} (max $0.62$ ever observed) and \textbf{CSD}
(max $0.68$) fail C2 --- replicating the audit's range finding and the coding
prior mismatch of Section~\ref{sec:livestudy}; \textbf{DBO}
(sd $0.037$, $r=0.76$ with TCA), \textbf{VDI} (sd $0.038$, $r$ with raw TAS
$=+0.00$) and \textbf{SHL} (sd $0.033$) fail C1. All five are
\emph{demoted, not deleted}: their detectors, per-step annotations, fix
generation (GAR still drives \code{goal\_anchor}) and the AVS verbosity alert
(VDI/SHL/CCR's actual job) run unchanged; they simply default to weight zero,
and a weights file can re-enable any of them. The nine surviving metrics keep
their previous weights, renormalised by the engine (non-zero sum $0.92$;
Table~\ref{tab:metrics}).

The corpus mean TAS rises from $73.5$ to $78.4$ under the pruned composite ---
mechanically expected, since two range-broken metrics had subtracted a
near-constant penalty from every trace while discriminating none. CCR is a
note of honesty in both directions: it \emph{passes} the pre-stated criteria
(sd $0.065$, max $0.99$) despite a near-zero correlation with the composite,
and it stays --- the criteria were chosen before the data was read, and they
cut both ways. The evaluation is regenerated by one command and is expected
to be re-run as the corpus grows; the criteria, not the current pass/fail
list, are the commitment.

\subsection{Adaptive sampling (preliminary)}
Table~\ref{tab:sampling} reports first-attempt pass rate and total tokens over
$50$ tau-bench airline tasks (gpt-4o), \textbf{single seed}. On point estimates,
self-consistency and adaptive-$K$ achieve the highest pass rate; adaptive-$K$
uses far fewer tokens than naive $K{=}5$. With $n=50$, the standard error on a
pass rate is roughly $7$ percentage points, so the gap between the top two
configurations is \emph{within noise}. We report this as motivating evidence,
not a significance result.

\begin{table}[t]
\centering
\caption{Sampling strategies on 50 tau-bench airline tasks, gpt-4o,
\textbf{single seed, no confidence intervals}. Differences within $\sim$1--2
standard errors should not be over-interpreted.}
\label{tab:sampling}
\begin{tabular}{lrr}
\toprule
Strategy & Pass@1 & Total tokens \\
\midrule
Single ($K{=}1$)        & 38\% & \phantom{0}63{,}248 \\
Naive $K{=}5$           & 40\% & 281{,}519 \\
Adaptive-$K$            & 46\% & 245{,}888 \\
Self-Consistency $K{=}5$ & 48\% & 136{,}942 \\
\bottomrule
\end{tabular}
\end{table}

\subsection{Performance}
Profiling with the production bag-of-words backend showed two dominant costs:
the similarity function re-tokenised both texts on every call ($9{,}534$ calls
covering only $191$ distinct texts on a 100-step trace, i.e.\ $\approx 99\%$
redundant tokenisation), and the context metric re-joined the whole prefix per
step ($O(n^2 \cdot \mathrm{len})$); together they were $486$\,ms of in-process
analysis at $500$ steps. The default similarity closure now memoises per-text
term-frequency vectors and CCE builds its prior-n-gram set incrementally with a
boundary-spanning tail; both changes are equivalence-tested to identical
output. Measured after the change, a $500$-step trace audits in $\approx
45$\,ms \emph{end-to-end} (process start included) and typical traces (tens of
steps) in low single-digit milliseconds. With the bounded look-back windows in
SRR and ISR, per-trace cost scales roughly linearly with step count; the HTTP
server additionally caps accepted traces at $50{,}000$ steps as a
denial-of-service guard.

\section{Threats to Validity and Limitations}
\label{sec:limits}
We consider this section the most important part of the report. TraceRazor is a
heuristic engineering tool; the following limit any scientific interpretation of
its outputs.

\paragraph{Construct validity of TAS.} TAS has not been validated against an
external criterion (human-labelled ``recoverable waste,'' or measured token
reduction at constant task success). It should be treated as an \emph{ordinal,
within-project} indicator, not a cross-agent or absolute efficiency percentage.

\paragraph{Uncalibrated weights and thresholds (by default).} The \emph{default}
composite weights (Table~\ref{tab:metrics}), the TVI constants in
Eq.~\eqref{eq:tvi}, the AVS coefficients, and per-metric thresholds (e.g. the
$0.65$ similarity cutoff) are author-chosen. The calibration procedure
(Section~\ref{sec:calib}) can fit them to data, but the real-data test on
tau-bench (Section~\ref{sec:realcalib}) found that no fit of the current metrics
predicts real recoverable waste (negative cross-validated $R^2$). So the shipped
score remains an ordinal within-trace heuristic by design, not a calibrated
predictor of cross-run savings, and the TVI/AVS constants and thresholds are
likewise uncalibrated. The validity audit in Section~\ref{sec:metric-eval}
quantifies the consequence on real data: realised metric influence scales with
empirical variance rather than the documented weights (TUR carried $28\%$ of
raw-TAS variance at a $7.9\%$ nominal weight under the fourteen-metric
composite), and the task-value multiplier dominates the final score
($r=0.89$). The rationalisation of Section~\ref{sec:rationalisation} removes
the metrics those findings indicted, but the surviving nine weights remain
heuristic --- pruning is not calibration --- and we have not performed a
sensitivity analysis showing conclusions are robust to perturbing the
constants.

\paragraph{Internally competing sub-metrics.} Some sub-metrics optimise in
opposite directions on the same similarity signal (redundancy penalises
inter-step similarity; continuity rewards it). There is consequently no single
joint optimum, and an agent can trade one against another.

\paragraph{Savings are estimates, not outcomes.} As stressed in
Section~\ref{sec:fixes}, reported token/cost savings are heuristic projections.
A measured saving requires applying the fixes, re-running the agent, and
comparing via \code{bench} at constant task quality. The one such measurement
performed to date (Section~\ref{sec:livestudy}) found the projection for the
\code{goal\_anchor} fix had the wrong sign on an on-track coding workload
(estimate accuracy $-102\%$), which should calibrate expectations for the
remaining, still-unmeasured fix types.

\paragraph{Scope of the measured study.} The intervention study covers one
agent harness, one model, one domain (small Python repairs), six tasks with
two replicates, and a single fix type --- the only one that fired. Its CIs are
honest about small-$n$ width and must not be extrapolated across models or
domains; the other safe fix types (verbosity, hedging, reformulation) remain
unmeasured.

\paragraph{Metric priors on coding agents.} The detection side carries a
conversational-agent prior that the coding corpus exposes: CSD scored every
coding trace as drifting (read$\to$edit$\to$test steps are legitimately
semantically distant), trajectory path entropy triggered \code{goal\_anchor}
on every trace, and TCA counts an agent's deliberate run of a failing test
suite --- an expected-failure probe, standard practice for coding agents ---
as a tool misfire. These penalties are identical across compared conditions,
so before/after deltas stand, but absolute scores on coding traces are
conservative-biased, and remediation text must therefore be cost-safe even
when triggered spuriously (the round-1 lesson of
Section~\ref{sec:livestudy}).

\paragraph{Closed-loop self-validation.} An optional adherence metric checks
whether the tool's own fixes improved the tool's own metrics. This is
self-validation, not external ground truth, and should not be read as evidence
that fixes improve real outcomes.

\paragraph{Sampling statistics.} The sampling benchmark is single-seed,
$n=50$, with no confidence intervals or significance tests. Multi-seed runs over
more tasks and domains are needed before any ranking of strategies is asserted.
We note that an earlier evaluation harness selected self-consistency answers
using benchmark gold outputs (an oracle); the \emph{shipped} selector uses honest
majority vote, and the oracle-based path should not be used for headline claims.

\paragraph{Substitutability circularity.} As described in
Section~\ref{sec:subst}, the shipped classifier evaluation has label leakage by
construction and is a smoke test only.

\paragraph{Similarity backend and language.} The default bag-of-words similarity
is shallow; verbosity word-lists are English-only. The compression-ratio proxy
uses a single global constant and can mis-fire on structured (e.g. JSON) content.

\section{Reproducibility}
\label{sec:repro}
The Rust workspace builds with \code{cargo build --release}; the Python package
installs with \code{pip install -e .}. Audits are reproduced with
\code{tracerazor audit <trace> --format json}; the real-trace benchmark with
\code{python benchmark/run\_benchmarks.py}; the metric-effectiveness
evaluation behind the composite (Section~\ref{sec:rationalisation}) with
\code{python -m benchmark.metric\_effectiveness}; the sampling data is shipped
as JSON; and the test suites run with \code{cargo test --workspace} (260 tests)
and \code{pytest} (263 tests). The audit table and benchmark numbers in
Section~\ref{sec:eval} were produced by the current binary on the committed
traces. The measured intervention study (Section~\ref{sec:livestudy}) ships
its complete dataset --- both rounds' traces, applied fixes, prompt patches
and run logs (cost, turns, session ids) --- under \code{benchmark/live/}, so
its measurement step,
\code{python -m benchmark.case\_study -{}-pairs-dir benchmark/live/traces},
re-runs without any API spend; re-running the agent halves requires a
Claude~Code install and cost $\approx\$1.30$ in total at the study's scale.

\section{Conclusion}
TraceRazor is a fast, interpretable, offline-first tool for finding and
explaining token waste in LLM agent traces, with optional sampling and
substitutability pillars. Its engineering is solid and its outputs are
actionable, but its central score is an uncalibrated ordinal heuristic and its
quantitative claims are either estimates or preliminary. We have tried to be
explicit about this throughout. Notably, when we calibrated on real tau-bench
runs (Section~\ref{sec:realcalib}) the metrics did not predict recoverable waste
(negative cross-validated $R^2$), which is the most important finding here: the
score is a useful within-trace diagnostic, not a validated predictor of cross-run
savings. Acting on that, we added observation-accumulation features
(Section~\ref{sec:betterfeatures}) that raise the cross-validated $R^2$ into the
positive range on two independent datasets ($+0.08$ and $+0.12$), the first
signal here that predicts real recoverable waste and replicates; the absolute
magnitude is still modest ($\approx 0.1$), and we found that obvious extensions
(path/length features, non-linear models) do not push past it
(Section~\ref{sec:betterfeatures}). The same evidence-gated discipline now
governs the composite itself: the validity audit's findings were re-derived
over 61 real traces and the five metrics they indicted (GAR, CSD, DBO, VDI,
SHL) were demoted to detection-only diagnostics
(Section~\ref{sec:rationalisation}). Of the four items we previously named as the path to a defensible
measurement, the second --- replace estimated savings with measured
before/after re-runs at constant task success --- is now delivered
(Section~\ref{sec:livestudy}), and its first finding was that our own
auto-applied fix cost tokens where the heuristic projected savings
($-5.6\%$ measured, estimate accuracy $-102\%$); rewriting the fix from
that evidence and re-measuring yielded a cost-neutral patch
($+0.7\%$, CI spanning zero) at constant $12/12$ pass rate. We consider the
willingness of the harness to indict its own product the strongest evidence
that the methodology is sound. The remaining path: (1) richer signal beyond
structural features (semantic/path-embedding features, per-token attribution)
plus calibration on the target's own production runs (the procedure of
Section~\ref{sec:calib} is in place and now has a feature set that beats the
mean); (2$'$) extend the measured loop to a workload where drift remediation
has something to recover (GAR failing), and to the still-unmeasured fix
types; (3) re-run the sampling study multi-seed with confidence intervals and
a deployable (non-oracle) selector; and (4) define and label substitutability
on real data. Until then, TraceRazor is best understood as a high-quality
heuristic and a reproducible harness --- now with a working, cheap, committed
procedure for catching its own mispriced advice --- which is how it is
documented and shipped.

\begin{thebibliography}{99}
\small
\bibitem{tale} M. Han et al. \emph{Token-Budget-Aware LLM Reasoning (TALE)}. ACL, 2025.
\bibitem{selfbudgeter} Q. Zhao et al. \emph{SelfBudgeter: Adaptive Token Allocation for Reasoning}. 2025.
\bibitem{stepsurvey} J. Lee et al. \emph{Evaluating Step-by-step Reasoning Traces: A Survey}. 2025.
\bibitem{dualformer} D. Su et al. \emph{Dualformer: Controllable Fast and Slow Thinking}. 2024.
\bibitem{steppruner} Y. Wu et al. \emph{Step Pruner: Efficient Reasoning in LLMs}. 2025.
\bibitem{efficientsurvey} S. Feng et al. \emph{Efficient Reasoning Models: A Survey}. 2025.
\bibitem{toolchain} Y. Pan et al. \emph{ToolChain*: Efficient Action-Space Search for LLM Agents}. NeurIPS, 2024.
\bibitem{budget} M. Hassid et al. \emph{Reasoning on a Budget}. 2025.
\bibitem{agenteval} A. Mohammadi et al. \emph{Evaluation and Benchmarking of LLM Agents}. KDD, 2025.
\bibitem{verbosity} W. Shi et al. \emph{Verbosity Bias in Preference Labeling / LLM Responses}. 2024.
\bibitem{selfconsistency} X. Wang et al. \emph{Self-Consistency Improves Chain-of-Thought Reasoning in Language Models}. ICLR, 2023.
\bibitem{taubench} S. Yao et al. \emph{tau-bench: A Benchmark for Tool-Agent-User Interaction in Real-World Domains}. Sierra Research, 2024. \url{https://github.com/sierra-research/tau-bench}
\bibitem{agentinstruct} A. Zeng et al. \emph{AgentTuning: Enabling Generalized Agent Abilities for LLMs}. arXiv:2310.12823, 2023. Dataset: \url{https://huggingface.co/datasets/zai-org/AgentInstruct}
\bibitem{trajhall} \emph{Trajectory-level hallucination taxonomies for tool-using agents}. arXiv:2605.24219, 2026; see also arXiv:2601.06818.
\bibitem{provagent} \emph{PROV-AGENT: Provenance graphs for agent workflows} (W3C-PROV extension). arXiv:2508.02866, 2025.
\bibitem{agentsmatter} S. Kapoor, B. Stroebl, Z. S. Siegel, N. Nadgir, A. Narayanan. \emph{AI Agents That Matter}. arXiv:2407.01502, 2024.
\bibitem{agentdiet} Y.-A. Xiao, P. Gao, C. Peng, Y. Xiong. \emph{Improving the Efficiency of LLM Agent Systems through Trajectory Reduction (AgentDiet)}. arXiv:2509.23586, 2025.
\bibitem{depo} S. Chen et al. \emph{DEPO: Dual-Efficiency Preference Optimization for LLM Agents}. arXiv:2511.15392, 2025.
\bibitem{spendmoney} L. Bai et al. \emph{How Do AI Agents Spend Your Money? Analyzing and Predicting Token Consumption in Agentic Coding Tasks}. arXiv:2604.22750, 2026.
\bibitem{tokenomics} M. Salim, J. Latendresse, S. Khatoonabadi, E. Shihab. \emph{Tokenomics: Quantifying Where Tokens Are Used in Agentic Software Engineering}. arXiv:2601.14470, 2026.
\end{thebibliography}

\end{document}
